% Shared preamble — included by all course main.tex files

% --- Layout ---
\usepackage[margin=1in]{geometry}

% --- Math ---
\usepackage{amsmath, amssymb, amsthm}

% --- Graphs ---
\usepackage{pgfplots}
\usepackage{tikz}
\usepackage{tikz-3dplot}
\pgfplotsset{compat=1.18}
\usetikzlibrary{calc}
\usetikzlibrary{decorations.pathreplacing}
\tikzset{
  number line/.style={<->, >=latex, semithick},
  sketch accent/.style={orange!75!black, thick},
  sketch point/.style={circle, fill=orange!75!black, inner sep=1.8pt}
}

% --- Boxes ---
\usepackage{tcolorbox}
\tcbuselibrary{theorems, breakable}

% --- Misc ---
\usepackage{enumitem}
\usepackage{hyperref}
\usepackage{pifont}

% --- Theorem-like environments ---

% Definition — green box, own counter
\newtcbtheorem[number within=section]{definition}{Definition}{%
  colback=red!3, colframe=red!40,
  fonttitle=\bfseries, breakable
}{def}

% --- Example box ---
\NewTColorBox{example}{ o }{%
  colback=gray!5,
  colframe=gray!60,
  fonttitle=\bfseries,
  title={Example\IfValueTF{#1}{: #1}{}},
  breakable
}

% Thinking — amber box for intuition, proof sketches, and tricks; keep sketches with their text
\NewTColorBox{thinking}{ o }{%
  colback=orange!5,
  colframe=orange!65!black,
  fonttitle=\bfseries,
  title={Thinking\IfValueTF{#1}{: #1}{}}
}

% Theorem — blue box, uses shared counter
\newtcbtheorem[number within=section]{theorem}{Theorem}{%
  colback=blue!3, colframe=blue!40,
  fonttitle=\bfseries, breakable
}{thm}

\newtcbtheorem[use counter from=theorem]{lemma}{Lemma}{%
  colback=blue!1, colframe=blue!20, coltitle=blue!50!black,
  fonttitle=\bfseries, breakable
}{lem}

\theoremstyle{plain}

\theoremstyle{remark}
\newtheorem*{remark}{Remark}
\newtheorem*{note}{Note}

% --- Proof layout ---
\definecolor{prooflabel}{gray}{0.30}
\newcommand{\UseUnindentedProofParagraphs}{%
  \AddToHook{env/proof/begin}{\setlength{\parindent}{0pt}}%
}
\newcommand{\inductionlabel}[1]{\textcolor{prooflabel}{\textit{#1}}}
\newcommand{\proofheadingbreak}{\leavevmode\par\nopagebreak}
\newlength{\proofgoalindent}
\setlength{\proofgoalindent}{1.5em}
\makeatletter
\newcommand{\proofgoal}[2]{%
  \par{\noindent
    % Keep an opening list label in its column, even when the item starts with a goal.
    \parshape 2 \@totalleftmargin \linewidth
    \dimexpr\@totalleftmargin+\proofgoalindent\relax
    \dimexpr\linewidth-\proofgoalindent\relax
    \hspace*{\proofgoalindent}\textcolor{prooflabel}{#1}\space#2\par}%
}
\makeatother
\SetEnumitemKey{proofparts}{font=\color{prooflabel}, align=left, labelindent=0pt, widest*=3,
  leftmargin=*, topsep=\smallskipamount, itemsep=\smallskipamount,
  parsep=0pt, partopsep=0pt}

% --- Common shortcuts ---
\newcommand{\R}{\mathbb{R}}
\newcommand{\N}{\mathbb{N}}
\newcommand{\Z}{\mathbb{Z}}
\newcommand{\Q}{\mathbb{Q}}
\newcommand{\C}{\mathbb{C}}
\newcommand{\ve}[1]{\langle #1 \rangle}   % vector brackets
\newcommand{\norm}[1]{\left\lVert #1 \right\rVert}
\newcommand{\abs}[1]{\left\lvert #1 \right\rvert}
\newcommand{\dd}[1]{\,\mathrm{d}#1}       % differential
\newcommand{\pd}[2]{\frac{\partial #1}{\partial #2}}  % partial derivative
\newcommand{\xmark}{\ding{55}}
